%% Springer Nature sn-jnl template. Compile in the official Springer Nature
%% LaTeX template (Overleaf: "Springer Nature LaTeX Template"), which provides
%% sn-jnl.cls and the .bst files. Reference style sn-basic = author-year, which
%% matches the Machine Learning journal. Put refs.bib and the figures in the same
%% folder. If citations error, confirm the class option matches the journal's
%% reference style; if a theorem env clashes, remove the matching \newtheorem below.
\documentclass[sn-basic]{sn-jnl}

\usepackage{graphicx}
\usepackage{amsmath,amssymb,amsfonts}
\usepackage{booktabs}

\theoremstyle{thmstyleone}
\newtheorem{proposition}{Proposition}
\newtheorem{assumption}{Assumption}
\theoremstyle{thmstyletwo}
\newtheorem{remark}{Remark}

\raggedbottom

\begin{document}

\title[Shift-Robust Selective Generation for RAG]{Knowing When Not to Answer:
Noise-Conditional Hallucination Risk Estimation and Shift-Robust Selective
Generation for Retrieval-Augmented Language Models}

\author*[1]{\fnm{Anmita} \sur{Das}}\email{aadas@student.unimelb.edu.au}
\author[1]{\fnm{Shoumya} \sur{Chowdhury}}\email{shoumyac@student.unimelb.edu.au}
\author[2]{\fnm{Sushanta} \sur{Paul}}\email{sushanta.customs@gmail.com}

\affil[1]{\orgdiv{School of Computing and Information Systems},
\orgname{University of Melbourne}, \orgaddress{\city{Melbourne}, \country{Australia}}}
\affil[2]{\orgname{National Board of Revenue (Bangladesh Customs)},
\orgaddress{\city{Dhaka}, \country{Bangladesh}}}

%% ORCID (enter in the submission system):
%% Anmita Das 0009-0000-1401-0494; Shoumya Chowdhury 0009-0005-5552-1094

\abstract{Retrieval-augmented generation (RAG) is widely deployed to reduce
hallucination in large language models, yet retrieval is itself a noise
source: irrelevant, contradictory, or misleading evidence can induce
hallucinations a model would not otherwise produce. We study selective
generation for RAG, where a system answers only when its hallucination
risk is statistically controlled, under the realistic condition that the
retrieval-noise distribution shifts between calibration and deployment.
Using a controlled two-regime testbed that separates parametric from
retrieval-induced error, we establish three behavioural regularities: a
single unopposed false document is absorbed at 85 to 100 percent across
three generators of two architectures, even against facts a model knows
(context dominance); contradictions at the first context position are
more than twice as damaging as at the last (misinformation primacy); and
under conflicting evidence all errors are confident assertions rather
than abstentions (silent failure). We then show that a noise-conditional
risk estimator using observable retrieval descriptors predicts
hallucination far better than model-internal uncertainty, beating
sampling-based semantic entropy by 0.31 to 0.40 AUROC across six
instruction-tuned model families. Standard conformal risk control and its
covariate-shift remedies all violate a 10 percent target in 200 of 200
resplits under noise shift; we introduce risk-binned Mondrian conformal
risk control, prove its validity for bin-measurable shifts, and show it
eliminates violations at competitive answer rates, with a label-free
deployment monitor that makes its assumption checkable. Results hold
across eight generators from 80M to 9B parameters, five model families,
two architectures, and two content domains. Code and data are released.}

\keywords{retrieval-augmented generation, hallucination, selective generation,
conformal risk control, distribution shift, uncertainty quantification}

\maketitle

\section{Introduction}\label{sec:intro}

Retrieval-augmented generation \citep{lewis2020rag}
grounds large language models (LLMs) in external corpora and is
routinely presented as a mitigation for hallucination. The
argument rests on an assumption that deployment rarely satisfies:
that retrieved evidence is relevant and correct. Production retrieval
pipelines return passages that are off-topic, partially relevant,
stale, about confusable entities, or factually wrong
\citep{chen2023rgb,wang2024astuterag,cuconasu2024power}. Worse, the
knowledge corpus itself can be corrupted: \citet{zou2025poisonedrag}
show that injecting a handful of crafted documents reliably steers RAG
answers, and defenses remain weak. A growing literature documents that
LLMs over-trust retrieved context even against their own parametric
knowledge \citep{xie2024adaptive,kortukov2024conflicts,xu2024survey}, so that RAG
can paradoxically suppress appropriate abstention because added context
raises confidence even when that context is wrong.
When evidence is bad, RAG does not merely fail to help; it manufactures
fluent, confident errors.

For high-stakes applications the question is therefore not only how to
make RAG more robust, but when a RAG system should decline to answer.
Selective prediction has a long history in classification
\citep{chow1970,geifman2017selective}; for generation, conformal
methods provide finite-sample guarantees on error rates among answered
queries \citep{angelopoulos2024crc,abbasi2024conformalabstention,
quach2024conformal,mohri2024factuality}. Existing conformal treatments
of RAG, however, either certify population-level risk
\citep{kang2024crag} or assume the calibration data is exchangeable
with deployment \citep{li2024traq,feng2025conformalrag}. Retrieval
noise breaks exchangeability in a specific and practically common way:
the mixture of noise regimes at deployment (which retriever, which
corpus, which adversary) differs from the mixture at calibration. A
system calibrated on clean development traffic and deployed against a
degraded or poisoned corpus faces exactly this shift.

This paper studies per-decision risk control for RAG under
retrieval-noise shift. We ask three questions. First, can hallucination
be predicted before answering from observable signals, and how much do
retrieval-side descriptors add over the model-internal uncertainty
measures that dominate current practice
\citep{kuhn2023semantic,farquhar2024nature,kossen2024sep}? Second, how
do noise type, severity, and position causally shape hallucination and
abstention, and how do these effects interact with parametric
knowledge? Third, when standard conformal calibration fails under noise
shift, what restores validity, and at what cost in answer rate?

Causal answers to the second question require intervening on evidence
while holding queries fixed, and require separating parametric error
(the model would have been wrong anyway) from retrieval-induced error
(the noise caused the mistake). Naturalistic benchmarks support
neither. We therefore construct a controlled two-regime testbed. A
synthetic world of 240 fictional scientists guarantees that no
pretrained model holds parametric knowledge of the queried facts, so
closedbook hallucination is 100\% by construction and every correct
answer is attributable to retrieval. Sixty well-known real facts
(country capitals and currencies) provide the complementary regime in
which mid-size models answer correctly without context 60\% to 90\% of
the time. Typed noise is injected into a BM25 pipeline at graded
severities and controlled context positions: partial and total
irrelevance, contradictions at first and last rank, unopposed
misinformation, confusable-entity distractors, and compound mixtures.
Every generation is logged with retrieval-side, model-side, and
closedbook-contrast features and labeled exactly against the closed
answer space. The released dataset comprises 6{,}230 generations on the main
domain across three generators spanning encoder-decoder and
decoder-only architectures, plus 2{,}970 generations on an independent
second domain.

Our contributions are as follows.

\begin{enumerate}
\item \textbf{Behavioural regularities of noisy RAG}
(\S\ref{sec:phenomena}). We document three phenomena. Context
dominance: one unopposed false document yields hallucination rates of
85\% to 100\% across all three generators, including on facts the
models answer correctly without context. Misinformation primacy: a
contradiction at rank 1 is more than twice as damaging as the same
document at rank 4 (53\% versus 23\% hallucination without parametric
support; 30\% versus 0\% with it). Universal silent failure: under
every conflict-noise type, 100\% of errors are confident assertions;
behavioural abstention occurs only under irrelevance noise. The
negative-rejection ability targeted by current instruction tuning
\citep{chen2023rgb} detects absent evidence but not
deceptive evidence.
\item \textbf{Noise-conditional risk estimation}
(\S\ref{sec:estimation}). Combining retrieval descriptors, model-internal signals, and
closedbook-contrast features raises hallucination-prediction AUROC from 0.788
to 0.915, with the margin over model-internal signals widening as noise becomes
subtler and concentrated in conflict noise (AUROC 0.60 to 0.72), the regime in
which models never warn.
\item \textbf{Shift-robust selective generation}
(\S\ref{sec:conformal}). Under clean-heavy calibration and noise-heavy
deployment, marginal conformal risk control and its standard covariate
shift remedies, learned-group calibration and density-ratio weighting
\citep{tibshirani2019covariate}, violate a 10\% risk target in every
one of 200 resplits. Risk-binned Mondrian conformal risk control,
which certifies score bins with Clopper-Pearson upper bounds, is proved
valid for bin-measurable shifts (Proposition~\ref{prop:riskbin}),
eliminates violations entirely, answers 2.3 times more queries than a
worst-group bound, and remains valid across target levels
$\alpha \in \{0.10, 0.15, 0.20\}$, bin counts $B \ge 10$, and all
tested shift severities.
\item \textbf{Transfer analysis} (\S\ref{sec:transfer}). Risk
estimators transfer across eight generators from 80M to 9B parameters
and five model families with AUROC 0.87 to 0.96, yet
transplanted thresholds realize 1.7 to 2.7 times the target risk.
Discrimination is substantially model-general; calibration is not. The binned certification scheme supports cheap local recalibration.
\end{enumerate}

\section{Related Work}\label{sec:related}

\subsection{Noise, conflict, and corruption in RAG}
Benchmarks and analyses have established that retrieval quality governs
RAG reliability. RGB isolates noise robustness and negative rejection
as core abilities and finds models lacking in both
\citep{chen2023rgb}. RAGTruth provides word-level hallucination
annotations for RAG outputs \citep{wu2024ragtruth}. RAAT types
retrieval noise and trains adapters against it \citep{fang2024raat};
Astute RAG consolidates internal and external knowledge at inference
\citep{wang2024astuterag}. \citet{cuconasu2024power} report the
counterintuitive finding that random irrelevant documents can improve
RAG accuracy, while related distractors hurt; our graded-irrelevance
results refine this picture (\S\ref{sec:discussion}). On the
adversarial side, PoisonedRAG demonstrates effective knowledge
corruption attacks with few injected documents
\citep{zou2025poisonedrag}. Position effects are documented for long
contexts by \citet{liu2024lost}; we show the analogous effect for
misinformation specifically. Knowledge-conflict studies characterise
when models prefer context over memory
\citep{xie2024adaptive,kortukov2024conflicts,xu2024survey}. This
literature measures and manipulates behaviour; it does not predict
per-query hallucination risk or provide decision-level guarantees.

\subsection{Uncertainty quantification for LLMs}
Semantic entropy estimates uncertainty over meanings rather than
strings \citep{kuhn2023semantic,farquhar2024nature}; probes recover it
cheaply from hidden states \citep{kossen2024sep}. Self-evaluation signals include P(True) \citep{kadavath2022language}. These methods treat the
generator as a closed system. None conditions on the quality or
structure of retrieved evidence, and our results quantify the cost of
that omission under retrieval noise.

\subsection{Selective prediction and conformal guarantees}
Selective classification with a reject option was introduced for deep networks by \citet{geifman2017selective}. Conformal risk control bounds the
expectation of monotone losses \citep{angelopoulos2024crc}; conformal abstention applies this to hallucination \citep{abbasi2024conformalabstention}. For RAG, TRAQ ensembles
conformal sets over retrieval and generation \citep{li2024traq};
Conformal-RAG filters sub-claims with group-conditional coverage
\citep{feng2025conformalrag}; C-RAG certifies an upper bound on
population-level generation risk, including under bounded distribution
drift \citep{kang2024crag}. Weighted conformal prediction corrects for
known covariate shift via likelihood ratios
\citep{tibshirani2019covariate}, adaptive conformal inference tracks
unknown drift online \citep{gibbs2021aci}, and the limits of
distribution-free conditional inference are themselves characterised by
\citet{barber2021limits}.
The guarantees in this literature are of three kinds, and none is the one a
selective-answering system needs: conformal risk control and C-RAG bound
population or marginal risk \citep{angelopoulos2024crc,kang2024crag}; conditional
conformal methods and Conformal-RAG bound conditional \emph{coverage} rather than
risk \citep{gibbs2025conditional,feng2025conformalrag}; and covariate-shift
methods, including weighted conformal risk control, restore only \emph{marginal}
guarantees under a shifted test distribution
\citep{tibshirani2019covariate,gibbs2021aci,zecchin2025wcrc}. Since exact
feature-conditional coverage is impossible distribution-free
\citep{barber2021limits}, we instead partition the feature space into a small
number of risk bins. We show empirically that likelihood-ratio weighting and
unsupervised grouping both fail under noise-regime selection, and that population
certificates and per-decision validity diverge sharply here; our risk-binned
Mondrian conformal risk control is, to our knowledge, the only method delivering
per-query, risk-bin-conditional control of the answer-or-abstain risk that stays
valid under measurable retrieval-noise covariate shift, needing neither ratio
estimation nor exchangeability, only that the calibrated score resolves the shift.

\section{Problem Setup and Method}\label{sec:method}

\subsection{Setting and observable features}
Let $q$ denote a query, $D = \{d_1, \dots, d_k\}$ the retrieved
context, $a = f_\theta(q, D)$ the generated answer, and
$y \in \{0, 1\}$ the indicator that $a$ is a hallucination, defined as
a confident assertion whose normalised form contradicts the gold
answer. Behavioural abstentions (the model answers ``unknown'' or
equivalent) form a separate outcome class and are analysed separately.
Each example carries a feature vector
$x = (x^{\mathrm{ret}}, x^{\mathrm{mod}}, x^{\Delta})$.

Retrieval-side descriptors $x^{\mathrm{ret}}$: mean, minimum, and top-1
BM25 relevance of the context; the top-1 to top-2 margin; score
dispersion; an inter-document conflict flag computed by extracting
values for the queried attribute from each context document and testing
for disagreement; the number of extracted claims; and a
context-presence indicator. All are computable at decision time without
labels.

Model-side signals $x^{\mathrm{mod}}$: mean and minimum token
log-probability of the generated answer, mean, maximum, and first-token
predictive entropy, and answer length, in the spirit of standard
white-box uncertainty measures.

Closedbook-contrast features $x^{\Delta}$: the change in mean answer
log-probability and entropy relative to the same query answered with no
context, and an indicator of whether the answer changed. These
operationalise the confidence-gain idea of \citet{bi2025ckplug} as risk
features. They require one additional closedbook pass per query, which
is cheap relative to retrieval.

\subsection{Risk estimation}
We estimate $\hat r(x) \approx \Pr[y = 1 \mid x]$ with
gradient-boosted trees (150 trees, depth 3), evaluated with
question-grouped five-fold cross-validation so that no question appears
in both training and test folds. Feature-family ablations isolate the
contribution of each signal class.

\subsection{Selective generation and risk-binned certification}
\label{sec:methodconformal}
A selection rule answers query $x$ iff $\hat r(x) \le \tau$. Conformal
risk control (CRC) selects $\tau$ on calibration data so that the
hallucination rate among answered queries is at most a target $\alpha$,
with finite-sample correction. The guarantee requires calibration and
deployment to be exchangeable. Retrieval-noise shift, modelled as a
change in the mixture of noise regimes, violates this.

We propose risk-binned Mondrian CRC. Partition the calibration score
range into $B$ quantile bins $\{I_b\}_{b=1}^{B}$. For each bin compute
the Clopper-Pearson upper confidence bound $\widehat p^{\,+}_b$ at
level $1 - \delta$ on the within-bin hallucination rate, and certify
bin $b$ iff $\widehat p^{\,+}_b \le \alpha$. At deployment, answer iff
the score falls in a certified bin. The scheme never extrapolates a
threshold into score regions whose risk it has not measured; it only
reuses bins whose risk it has certified.

\begin{assumption}[Bin-measurable covariate shift]\label{ass:shift}
The deployment distribution $Q$ and calibration distribution $P$
satisfy (i) $Q(y \mid x) = P(y \mid x)$, and (ii) the likelihood ratio
$dQ/dP(x)$ is constant on each bin $I_b$.
\end{assumption}

Condition (ii) formalises the requirement that the score resolves the
shift: deployment may reweight bins arbitrarily, including
adversarially toward high-risk bins, but may not redistribute mass
within a bin. Selection over noise regimes that the risk score
separates is of exactly this form.

\begin{proposition}[Validity under bin-measurable shift]
\label{prop:riskbin}
Under Assumption~\ref{ass:shift}, with probability at least
$1 - B\delta$ over the calibration draw, the hallucination rate among
answered deployment queries is at most $\alpha$, simultaneously for
every deployment distribution $Q$ satisfying the assumption.
\end{proposition}

\begin{proof}
Fix a bin $I_b$. By Assumption~\ref{ass:shift}(ii) the conditional
distribution of $x$ given $x \in I_b$ coincides under $P$ and $Q$;
with (i), the within-bin risk $p_b = \Pr[y = 1 \mid x \in I_b]$ is
distribution-independent. Calibration labels within $I_b$ are i.i.d.\
Bernoulli($p_b$), so $\Pr[p_b \le \widehat p^{\,+}_b] \ge 1 - \delta$
by the Clopper-Pearson construction. A union bound over the $B$ bins
gives, with probability at least $1 - B\delta$, $p_b \le \alpha$ for
every certified bin. The deployment risk among answered queries is a
$Q$-weighted convex combination of certified-bin risks and is therefore
at most $\alpha$ for any bin reweighting, which covers every $Q$ in the
assumption class simultaneously.
\end{proof}

\begin{remark}
The guarantee requires no estimate of the likelihood ratio. This is the
step at which weighted conformal methods fail in our experiments: when
noise regimes overlap in feature space, the estimated ratio
under-corrects precisely in the high-risk tail. The price of
certification is conservativeness, quantified empirically in
\S\ref{sec:conformal} and \S\ref{sec:sensitivity}: within-bin score
ordering is discarded, and small bins certify only with generous
margins. The union-bound factor $B\delta$ ties the bin count to the
confidence level; we use $B = 10$ and $\delta = 0.10$ in the main
experiments and sweep both in \S\ref{sec:sensitivity}.
\end{remark}

We next characterise the price of certification in answer rate, which
explains when the method is useful and how to choose the bin count.

\begin{proposition}[Answer rate and its decomposition]\label{prop:eff}
Let $\pi_b = \Pr_Q[\hat r(x) \in I_b]$ be the deployment mass of bin
$I_b$ and let $C \subseteq \{1, \dots, B\}$ index the certified bins.
The answer rate of risk-binned Mondrian CRC is
$\rho = \sum_{b \in C} \pi_b$, and a bin is certified whenever its
empirical risk satisfies
$\widehat p_b \le \alpha - \varepsilon_b(\delta, n_b)$, where
$\varepsilon_b(\delta, n_b)$ is the Clopper-Pearson half-width at
calibration count $n_b$. Consequently
\[
\rho \;\ge\; \sum_{b:\, p_b \le \alpha - 2\varepsilon_b} \pi_b ,
\]
so the achievable answer rate is governed by the deployment mass lying
in score bins whose true risk is bounded away from $\alpha$ by twice
the finite-sample margin.
\end{proposition}

\begin{proof}
The answer-rate identity is immediate since answered queries are exactly
those whose score lands in a certified bin. For the bound, the empirical
within-bin risk concentrates around $p_b$ within
$\varepsilon_b(\delta, n_b)$ with probability $1 - \delta$ by the
Clopper-Pearson interval; whenever $p_b \le \alpha - 2\varepsilon_b$,
the upper bound $\widehat p^{\,+}_b = \widehat p_b + \varepsilon_b
\le p_b + 2\varepsilon_b \le \alpha$, so the bin is certified. Summing
the deployment mass of these bins lower-bounds $\rho$.
\end{proof}

Proposition~\ref{prop:eff} yields two practical consequences. First,
the method answers nothing only when essentially all deployment mass
sits in bins whose true risk exceeds $\alpha$, that is, when the
deployment stream is genuinely almost all high-risk; in that regime
abstention is the correct behaviour. Second, because
$\varepsilon_b = O(\sqrt{p_b(1 - p_b)/n_b})$, increasing the bin count
$B$ shrinks each $n_b$ and widens $\varepsilon_b$, which lowers answer
rate, while decreasing $B$ coarsens the partition and risks violating
Assumption~\ref{ass:shift}(ii). The two effects bracket an interior
optimum, consistent with the bin-count sweep in
\S\ref{sec:sensitivity}, where $B = 10$ maximises answer rate subject
to validity.

\begin{figure}[tbp]
\centering
\fbox{\begin{minipage}{0.9\linewidth}
\textbf{Algorithm 1: Risk-binned Mondrian CRC for selective RAG}
\par\vspace{3pt}\hrule\vspace{4pt}
\textbf{Input:} calibration set $\{(x_i, y_i)\}_{i=1}^{n}$, risk scorer
$\hat r$, target $\alpha$, bins $B$, level $\delta$.
\begin{enumerate}\setlength{\itemsep}{1pt}
\item Compute scores $s_i = \hat r(x_i)$ and bin edges
$\{q_b\}_{b=0}^{B}$ as $B{+}1$ empirical quantiles of $\{s_i\}$.
\item For each bin $b = 1, \dots, B$: let
$S_b = \{i : q_{b-1} < s_i \le q_b\}$,
$k_b = \sum_{i \in S_b} y_i$, $n_b = |S_b|$; compute the
Clopper-Pearson upper bound
$\widehat p^{\,+}_b = \mathrm{BetaInv}(1 - \delta,\, k_b + 1,\, n_b - k_b)$;
certify bin $b$ if $\widehat p^{\,+}_b \le \alpha$.
\item At deployment, for query $x$: compute $s = \hat r(x)$, locate its
bin $b(s)$; answer if $b(s)$ is certified, otherwise abstain.
\end{enumerate}
\end{minipage}}
\caption{Risk-binned Mondrian conformal risk control.}
\label{alg:rb}
\end{figure}

\section{Experimental Design}\label{sec:design}

\subsection{Two-regime fact world}
The corpus contains 2{,}200 documents. For each of 240 fictional
scientists we generate five gold fact sentences (birth year, birth
city, field, invention, institution) from templates with filler
clauses; for each of 60 real facts (30 capitals, 30 currencies) we
write one gold sentence. For every question we add one confident
contradiction document asserting a plausible false value, framed as a
revised archival review or administrative reform; for every synthetic
question, one confusable-entity distractor whose subject differs from
the queried entity by one character; and 400 irrelevant topical
fillers. Synthetic entities make parametric knowledge impossible, so
closedbook hallucination is 100\% by construction; real facts provide
the regime where parametric knowledge exists (closedbook accuracy:
flan-t5-base 60\%, Qwen2.5-0.5B-Instruct 90\%).

\subsection{Conditions}
Eleven conditions per question. Clean retrieval; graded irrelevance
irr25, irr50, irr75, irr100 (one to four of the $k = 4$ context slots
replaced by topically plausible irrelevant documents, with the gold
document removed at 100\%); contra\_r1 and contra\_r4 (gold present,
contradiction at first or last rank); contra\_only (contradiction
present, gold absent); distractor; mixed (gold plus contradiction plus
irrelevant); closedbook.

\subsection{Models, pipeline, and labels}
BM25 retrieval over the full corpus; greedy decoding with at most 12 to
14 new tokens. Generators: flan-t5-base (250M, encoder-decoder, all 11
conditions, 3{,}300 generations), flan-t5-small (80M, 6 conditions,
1{,}800 generations), and Qwen2.5-0.5B-Instruct (decoder-only,
chat-templated prompts, 4 conditions on a 100-question subset, 410
generations), plus a 720-generation preliminary campaign used only for
the noise-subtlety comparison. Answers are normalised and matched
exactly against the closed answer space, with a pattern-based
abstention class. Labels are therefore programmatic and exact, which
removes LLM-judge circularity from the measurement loop. The harness
checkpoints atomically, resumes from interruption, and runs on
commodity CPU hardware; a companion notebook scales the identical
protocol to 7B-class models on GPU.

\subsection{Evaluation protocol}
Risk estimation uses question-grouped five-fold cross-validation;
paired bootstrap over questions (2{,}000 resamples) gives confidence
intervals for AUROC differences. Conformal experiments use 200
question-level resplits; in each, calibration is drawn clean-heavy
(noisy-condition rows subsampled to 30\%) and deployment noise-heavy
(clean rows subsampled to 50\%), emulating calibration on development
traffic followed by deployment against a degraded corpus. Sensitivity
analyses vary the target $\alpha$, the bin count $B$, and the
calibration noise retention from 0.9 down to 0.1.

\section{Results}\label{sec:results}

\subsection{Behavioural regularities of noisy RAG}\label{sec:phenomena}

\begin{table}[tbp]
\centering\small
\caption{Hallucination (H) and abstention (A) rates by condition and
regime, flan-t5-base ($n = 240$ synthetic and $60$ real per cell).
C denotes correct.}
\label{tab:dose}
\begin{tabular}{lcccc}
\toprule
& \multicolumn{2}{c}{Synthetic (no param.\ knowledge)} &
\multicolumn{2}{c}{Real (param.\ knowledge)} \\
\cmidrule(lr){2-3}\cmidrule(lr){4-5}
Condition & H & A & H & A \\
\midrule
clean        & 0.00 & 0.00 & 0.00 & 0.00 \\
irr25        & 0.03 & 0.00 & 0.00 & 0.00 \\
irr50        & 0.09 & 0.00 & 0.00 & 0.00 \\
irr75        & 0.16 & 0.00 & 0.00 & 0.00 \\
irr100       & 0.62 & 0.38 & 0.20 & 0.78 \\
contra\_r1   & 0.53 & 0.00 & 0.30 & 0.00 \\
contra\_r4   & 0.23 & 0.00 & 0.00 & 0.00 \\
contra\_only & 1.00 & 0.00 & 1.00 & 0.00 \\
distractor   & 0.17 & 0.00 & 0.00 & 0.00 \\
mixed        & 0.45 & 0.00 & 0.15 & 0.00 \\
closedbook   & 1.00 & 0.00 & 0.40 & 0.00 \\
\bottomrule
\end{tabular}
\end{table}



\paragraph{Context dominance.}
Under contra\_only, hallucination is 100\% in both regimes for flan-t5-base
(Table~\ref{tab:dose}). The decoder-only replication is equally stark:
Qwen2.5-0.5B-Instruct, which answers the real-fact questions correctly 90\% of
the time with no context, hallucinates on 85\% of them when a single unopposed
false document is present and on 100\% of the synthetic questions. Parametric knowledge arbitrates conflicts only when
supporting evidence is also retrieved; in its absence, context wins almost
absolutely. This sharpens the parametric-bias picture of
\citet{kortukov2024conflicts} and gives a security reading consistent with
\citet{zou2025poisonedrag}: a corpus-poisoning attack need not outvote the
truth, only keep it absent from the retrieved window, and we quantify the
per-document absorption rate that makes such attacks efficient.

\paragraph{Misinformation primacy.}
With gold evidence present, moving the contradiction from rank 1 to
rank 4 cuts hallucination from 0.53 to 0.23 without parametric support
and from 0.30 to 0.00 with it. Position in the context window is a
deployment-observable risk factor, complementary to the attention
position effects of \citet{liu2024lost}, and no existing uncertainty
or selective-RAG method conditions on it.

\paragraph{Universal silent failure.}
Across every conflict condition (contra\_r1, contra\_r4, contra\_only,
distractor, mixed), 100\% of errors are confident assertions; the model
abstained in 0 of those error cases. Under pure irrelevance the model
abstains on 38\% (synthetic) to 78\% (real) of affected queries. The
same asymmetry holds exactly for the decoder-only model. Training for
negative rejection \citep{chen2023rgb} therefore
addresses the loud failure mode and leaves the silent one untouched,
which is the central argument for external, evidence-aware risk
control.



\paragraph{Chat tuning increases conflict vulnerability.}
With gold and contradiction both at rank 1, the chat-tuned decoder hallucinates
on 73\% of real-fact queries against 30\% for flan-t5-base, despite higher
closedbook accuracy (90\% versus 60\%); stronger context obedience, an explicit
objective of RAG instruction tuning, appears to trade against conflict
robustness, a trade-off worth testing at larger scales with the released harness.

\subsection{Noise-conditional risk estimation}\label{sec:estimation}

\begin{table}[tbp]
\centering\small
\caption{Hallucination-risk prediction, question-grouped five-fold CV,
flan-t5-base ($n = 3{,}161$ answered; base rate 0.363).}
\label{tab:auroc}
\begin{tabular}{lc}
\toprule
Feature family & AUROC \\
\midrule
Retrieval-side descriptors        & 0.897 \\
Model-internal (logprob, entropy) & 0.788 \\
Closedbook contrast               & 0.700 \\
Model + contrast                  & 0.811 \\
Combined (noise-conditional)      & \textbf{0.915} \\
\bottomrule
\end{tabular}
\end{table}

The combined estimator improves on model-internal signals by 0.127
AUROC (95\% CI 0.115 to 0.140); the closedbook-contrast family alone
adds 0.023 (CI 0.013 to 0.035) over model-internal signals
(Table~\ref{tab:auroc}). The comparison with our preliminary campaign
is instructive. With crude noise (only total irrelevance and rank-1
contradictions), the combined-versus-model gap was 0.026; with the
graded and positional noise grid it is 0.127, because model-internal
discrimination falls from 0.905 to 0.788 while retrieval descriptors
hold (0.899 to 0.897). Subtle noise is precisely where context-agnostic
uncertainty fails and noise conditioning pays.

\paragraph{Where prediction remains hard.}
Per-condition diagnostics (out-of-fold scores) show near-perfect
discrimination on irrelevance noise (irr25 to irr100: AUROC 0.98 to
1.00) and substantially lower discrimination under conflict
(contra\_r1: 0.605; mixed: 0.596; contra\_r4: 0.722; distractor:
0.839). The estimator is well calibrated overall (ECE 0.024, 10
quantile bins), and feature importances are dominated by mean retrieval
relevance (0.45) followed by minimum token log-probability and the
relevance margin. The hard residual case, conflict-induced
hallucination, is exactly the silent regime of
\S\ref{sec:phenomena}; this conjunction, hardest to predict and never
self-signalled, is what the conformal layer must cover.

\subsection{Comparison with uncertainty-quantification baselines}
\label{sec:baselines}

\begin{table}[tbp]
\centering\small
\caption{Hallucination prediction against real UQ baselines computed on the
same generator (flan-t5-small, 120-question subset, five samples per query
for semantic entropy). Higher is better. The noise-conditional estimator
uses only label-free retrieval descriptors.}
\label{tab:baselines}
\begin{tabular}{lc}
\toprule
Signal & AUROC \\
\midrule
Semantic entropy \citep{kuhn2023semantic,farquhar2024nature} & 0.675 \\
$1 - $ P(True) \citep{kadavath2022language}                      & 0.710 \\
Noise-conditional estimator (ours, retrieval only)           & 0.893 \\
Noise-conditional estimator $+$ semantic entropy $+$ P(True) & \textbf{0.915} \\
\bottomrule
\end{tabular}
\end{table}

A natural question is whether the gains over model-internal log-probability
and entropy (\S\ref{sec:estimation}) survive comparison with purpose-built
uncertainty-quantification methods. We computed two standard baselines on
the same generator: discrete semantic entropy
\citep{kuhn2023semantic,farquhar2024nature} from five temperature samples,
using exact normalised-string equivalence as the semantic-equivalence
relation (which is exact for the closed answer space rather than approximated
by an entailment model), and P(True) self-evaluation
\citep{kadavath2022language}. Table~\ref{tab:baselines} shows the noise-conditional
estimator outperforms semantic entropy by 0.218 AUROC
(95\% CI 0.193 to 0.243) and P(True) by 0.183. Adding both signals to our
estimator yields only a further 0.022 (95\% CI 0.012 to 0.033), so the
predictive power resides in the observable retrieval structure, not in
sampling-based self-consistency. This is consistent with the
mechanism we isolate in \S\ref{sec:decomp}: sampling-based entropy tracks the
model's internal disagreement, which is informative when evidence is missing
but near blind when the model confidently absorbs a coherent falsehood, so under
retrieval noise the decisive signal is the quality of the evidence, which
model-internal uncertainty does not observe. The gap widens with model capability: on Qwen2.5-7B-Instruct, sampling-based
semantic entropy is barely better than chance at predicting
retrieval-induced hallucination (AUROC 0.536), while the noise-conditional
estimator reaches 0.892, a difference of 0.356 AUROC (95\% CI 0.340 to
0.373) -- larger than the 0.218 gap on the 80M model, the opposite of what a
small-model artefact would predict.

\subsection{Selective generation under noise shift}\label{sec:conformal}

\begin{table}[tbp]
\centering\small
\caption{Selective generation at target risk $\alpha = 0.10$ under
clean-heavy calibration and noise-heavy deployment; 200 resplits,
flan-t5-base. Violation is the fraction of resplits whose realized risk
exceeds $\alpha$.}
\label{tab:crc}
\begin{tabular}{lccc}
\toprule
Method & Realized risk & Violation & Answer rate \\
\midrule
Marginal CRC                    & 0.215 & 100\% & 70.5\% \\
Heuristic-group CRC             & 0.189 & 100\% & 60.9\% \\
Learned-group CRC ($k$-means)   & 0.213 & 100\% & 67.0\% \\
Density-ratio weighted CRC      & 0.176 & 100\% & 61.5\% \\
Worst-group CRC                 & 0.035 & 22\%  & 12.7\% \\
Risk-binned Mondrian CRC (ours) & \textbf{0.063} & \textbf{0\%} & \textbf{29.2\%} \\
\bottomrule
\end{tabular}
\end{table}



Table~\ref{tab:crc} contains the central result. Marginal CRC roughly doubles its target risk. The two textbook
covariate-shift remedies fail in every resplit as well: unsupervised
$k$-means groups in descriptor space do not align with risk strata, and
the logistic-regression density ratio under-corrects in the high-risk
tail because noise regimes overlap in feature space. The worst-group
bound is valid in most resplits but answers only 12.7\% of queries.
Risk-binned certification eliminates violations entirely at a 29.2\%
answer rate. The mechanism is the one identified in
\S\ref{sec:methodconformal}: certification reuses only measured risk,
so reweighting bins cannot break it, whereas every thresholding method
implicitly extrapolates.

\subsection{Sensitivity and robustness analyses}\label{sec:sensitivity}
Three sweeps (60 resplits per cell; full figures in the repository) confirm the
guarantee is not an artefact of one operating point. Across targets, risk-binned
certification is valid at $\alpha \in \{0.10, 0.15, 0.20\}$ with zero violations
and shows $8\%$ violations at the strict $\alpha = 0.05$ (consistent with the
certification level $\delta = 0.10$), while marginal CRC violates at every target
with realized risk roughly $2\alpha$. Across bin counts, $B = 5$ is too coarse
($25\%$ violations) while $B \in \{10, 20, 40\}$ are all valid, trading answer
rate ($0.29$ down to $0.20$) against granularity exactly as the union-bound
analysis predicts. Across shift severities, marginal risk grows monotonically
from $0.143$ to $0.291$ as calibration sees less noise, while risk-binned
certification stays controlled throughout.

\subsection{Robustness to within-bin shift and a deployment monitor}
\label{sec:selfdiag}
Proposition~\ref{prop:riskbin} assumes the shift is bin-measurable
(Assumption~\ref{ass:shift}(ii)); its one failure mode is a within-bin shift
that raises a certified bin's true risk without moving scores. Holding the
inter-document conflict flag out of the risk score so that it acts as a latent
within-bin risk driver, we construct an adversarial deployment that upweights
high-conflict items twelvefold. Even then the method degrades only mildly:
realized risk rises from 0.066 to 0.072 and the violation rate from $0\%$ to
$4\%$ against the 0.10 target, because the score already absorbs most of the
risk-relevant covariate structure. When stricter safety is required, the
held-out covariate can be monitored at deployment without labels: decertifying
a bin when its deployment conflict rate exceeds calibration flags 4.67 bins on
average under the adversarial within-bin shift but only 1.39 under benign
reweighting (a 3.4-fold contrast), making the method's one assumption checkable
in deployment rather than merely assumed.

\subsection{Generality across models, estimators, and domains}
\label{sec:generality}

\begin{table}[tbp]
\centering\small
\caption{Generality of the central result across eight generators
spanning 80M to 9B parameters and five model families (FLAN-T5, Qwen,
Microsoft Phi, Mistral, Google Gemma), two architectures
(encoder-decoder and decoder-only), and an independent second domain.
For every generator the combined risk estimator discriminates
hallucination well (AUROC 0.87 to 0.96), the sampling-based semantic
entropy baseline stays near chance (AUROC 0.52 to 0.64), marginal CRC
violates the $\alpha = 0.10$ target in essentially every resplit, and
risk-binned Mondrian CRC restores the guarantee (0\% to 4\%
violations). The six decoder-only instruction-tuned models each ran the
full eleven-condition grid (about 3{,}300 generations) with white-box
logprob and entropy features and bidirectional-NLI semantic entropy.
AUROC columns use question-grouped cross-validation.}
\label{tab:generality}
\begin{tabular}{llccccc}
\toprule
& & \multicolumn{2}{c}{Estimator AUROC} &
\multicolumn{1}{c}{Marg.} & \multicolumn{2}{c}{Risk-binned CRC} \\
\cmidrule(lr){3-4}\cmidrule(lr){5-5}\cmidrule(lr){6-7}
Generator & Domain & Comb. & SE & Viol. & Viol. & Ans.\ rate \\
\midrule
flan-t5-small (80M, enc-dec)    & scientists & 0.959 & ---   & 100\% & 0\% & 17.5\% \\
flan-t5-base (250M, enc-dec)    & scientists & 0.915 & ---   & 100\% & 0\% & 28.8\% \\
Qwen2.5-1.5B (dec-only)         & scientists & 0.871 & 0.535 & 100\% & 0\% & 35.3\% \\
Qwen2.5-3B (dec-only)           & scientists & 0.874 & 0.530 & 100\% & 4\% & 29.7\% \\
Phi-3.5-mini (3.8B, dec-only)   & scientists & 0.948 & 0.640 & 100\% & 0\% & 28.9\% \\
Mistral-7B-v0.3 (dec-only)      & scientists & 0.950 & 0.617 & 100\% & 0\% & 36.9\% \\
\textbf{Qwen2.5-7B-Instruct (dec-only)} & scientists & \textbf{0.892} & \textbf{0.536} & \textbf{100\%} & \textbf{1\%} & \textbf{36.7\%} \\
Gemma-2-9B-it (dec-only)        & scientists & 0.923 & 0.523 & 100\% & 1\% & 39.0\% \\
flan-t5-small (80M, enc-dec)    & companies  & 0.920 & ---   & 100\% & 0\% & 21.7\% \\
\bottomrule
\end{tabular}
\end{table}



\paragraph{Six instruction-tuned families.}
To rule out that the result is an artifact of one model or one
architecture, we ran the full eleven-condition grid on six decoder-only
instruction-tuned generators spanning 1.5B to 9B parameters and four
independent pretraining families: Qwen2.5 (1.5B, 3B, 7B), Microsoft
Phi-3.5-mini (3.8B), Mistral-7B-v0.3, and Google Gemma-2-9B. Each run
logs true token logprobs, predictive entropy, and bidirectional-NLI
semantic entropy, for about 3{,}300 generations per model. The combined
noise-conditional estimator discriminates hallucination with AUROC
between 0.871 and 0.950 on every family, while sampling-based semantic
entropy never rises above 0.640 and sits near chance (0.52 to 0.54) on
four of the six. The combined-minus-semantic-entropy advantage ranges
from $+0.308$ to $+0.400$ AUROC, and every per-model bootstrap interval
excludes zero by a wide margin, so the baseline gap is not specific to
any one generator. Marginal CRC violates the $\alpha = 0.10$ target in
100\% of resplits on all six models, and risk-binned Mondrian CRC
restores the guarantee on all six (realized risk 0.022 to 0.060,
violation 0\% to 4\%) at answer rates of 29\% to 39\%. The largest
semantic-entropy gap ($+0.400$) and the highest answer rate (39\%) both
occur on the largest model, Gemma-2-9B, so neither the estimator
advantage nor the efficiency of the guarantee degrades with scale.

\paragraph{Second domain.}
On an independent domain with disjoint subject matter (240 fictional technology
companies and 30 real chemical-element facts, 2{,}970 generations with
flan-t5-small), every phenomenon reproduces: unopposed misinformation is
absorbed at 100\% (synthetic) and 73\% (real), the rank-1 contradiction is 0.72
against 0.18 at rank 4, and silent failure is near-universal. The estimator
reaches AUROC 0.920, marginal CRC violates the target in every resplit, and
risk-binned CRC never violates it (realized risk 0.016). The regularities and
the guarantee are properties of noisy RAG, not of one corpus.

\paragraph{Scaling and the position effect.}
At 7B (Qwen2.5-7B-Instruct), negative rejection improves with scale, so silent
failure under pure irrelevance falls to zero, yet under unopposed misinformation
97\% of errors remain confident assertions. The position effect inverts relative
to the encoder-decoder models: a contradiction at the last context position is
more damaging than at the first (rank-1 0.38 versus rank-4 0.67), so position
remains a major observable risk factor whose damaging end is
architecture-dependent.

\paragraph{Estimator architecture.}
The signal is not tied to gradient boosting: on flan-t5-base, logistic
regression, random forest, boosting, and a small multilayer perceptron reach
AUROC 0.870, 0.922, 0.915, and 0.905, so the gain comes from the features rather
than the model class.

\paragraph{Recalibration.}
Isotonic recalibration of the already well-calibrated score (ECE 0.024) does not
improve the answer rate, so we report the method without it; on a poorly
calibrated base scorer it would be expected to help.

\subsection{Cross-model transfer}\label{sec:transfer}
Training the risk estimator on one generator and deploying it on the
other preserves discrimination but not calibration. From flan-t5-small
to flan-t5-base, AUROC is 0.827 and a source-calibrated threshold at
$\alpha = 0.10$ realizes risk 0.170; in the reverse direction, AUROC
is 0.908 and realized risk is 0.269. Rankings of risky inputs are
substantially model-general; absolute risk levels are not. Portable
selective RAG therefore requires local recalibration, which binned
certification supports at low cost since bins can be re-certified on a
small labelled sample from the target system. We register the
prediction that the same dissociation governs cross-lingual transfer,
for instance from English to Bengali retrieval corpora, and leave its
test to follow-up work on naturalistic multilingual data.

\subsection{Naturalistic validation on real retrieval}
\label{sec:naturalistic}
The controlled world isolates each noise mechanism, but a referee may
reasonably ask whether the effects survive genuine retrieval. We
therefore rebuilt the pipeline on SQuAD: a 300-paragraph pool with one
question per paragraph, real BM25 and dense (e5 plus FAISS) retrieval,
the same typed-noise grid, and PoisonedRAG poisoning
\citep{zou2025poisonedrag} as an adversarial condition, run on two
instruction-tuned generators (Qwen2.5-7B and Mistral-7B-v0.3) with
white-box features and discrete semantic entropy.
Table~\ref{tab:natauroc} reports the estimator on the answered subset.
The combined estimator reaches AUROC 0.84 to 0.86 on genuine retrieval
for both generators, and BM25 and dense retrieval give almost identical
numbers (0.842 versus 0.840 for Qwen, 0.852 versus 0.857 for Mistral),
so the signal is retriever-independent. The advantage over semantic
entropy is positive and significant in every cell ($+0.14$ to $+0.31$
AUROC, all confidence intervals excluding zero), though smaller than in
the controlled world because Mistral's semantic entropy is itself more
informative here (0.71).

\begin{table}[tbp]
\centering\small
\caption{Naturalistic validation on real SQuAD retrieval. Combined risk
estimator versus the semantic entropy (SE) baseline on the answered
subset, for two generators and two retrievers. Combined-minus-SE shows a
bootstrap 95\% confidence interval.}
\label{tab:natauroc}
\begin{tabular}{llccc}
\toprule
Generator & Retriever & Combined & SE & Combined $-$ SE [95\% CI] \\
\midrule
Qwen2.5-7B & BM25  & 0.842 & 0.533 & $+0.309$ [$+0.284, +0.334$] \\
Qwen2.5-7B & dense & 0.840 & 0.545 & $+0.295$ [$+0.267, +0.325$] \\
Mistral-7B & BM25  & 0.852 & 0.714 & $+0.138$ [$+0.117, +0.159$] \\
Mistral-7B & dense & 0.857 & 0.693 & $+0.164$ [$+0.141, +0.186$] \\
\bottomrule
\end{tabular}
\end{table}

Per-condition hallucination on real retrieval confirms the controlled ordering: PoisonedRAG poisoning is the strongest
inducer (0.74 to 0.85 across models), while the clean condition already
hallucinates at 0.15 to 0.28. Because the clean base rate is this high, a
0.10 selective-risk target sits below it for these models, so the
realistic operating point is $\alpha = 0.20$. At that target the
dissociation reproduces directionally: marginal CRC violates between 47\%
and 62\% of resplits, whereas risk-binned Mondrian CRC holds (0\% to 1\%
violations, realized risk 0.09 to 0.12) at a 21\% to 23\% answer rate. We
frame this honestly. The controlled world produces the dramatic
headline that marginal CRC violates in essentially every resplit;
naturalistic noise is milder and the base rates higher, so marginal CRC
violates about half the resplits while the binned fix restores validity
at a usable answer rate. The naturalistic run corroborates the direction
of the controlled result on real data rather than reproducing its
magnitude.



\subsection{Why semantic entropy fails: missing-evidence versus
misinformation}\label{sec:decomp}
The baseline collapse and the naturalistic milder gap have a single
explanation, which also reconciles the controlled and natural results.
Within the injected naturalistic conditions we split failures into two
mechanisms and measure each signal's AUROC for predicting hallucination
(Table~\ref{tab:decomp}). \emph{Missing-evidence} failures (closed-book
and graded irrelevance) leave the model with no usable passage;
\emph{misinformation} failures (rank-1 contradiction, contradiction-only,
poisoning, and mixed) place a coherent but false passage in context.

\begin{table}[tbp]
\centering\small
\caption{Decomposition of naturalistic injected failures into
missing-evidence and misinformation, with each signal's AUROC for
predicting hallucination. Semantic entropy (SE) is informative on
missing-evidence but near chance on misinformation; retrieval
descriptors are the mirror image; the combined estimator captures both.}
\label{tab:decomp}
\begin{tabular}{lllcccc}
\toprule
Generator & Failure mode & $n$ & H-rate & SE & Retrieval & Combined \\
\midrule
Qwen2.5-7B & missing-evidence & 51  & 0.69 & 0.668 & 0.513 & 0.420 \\
Qwen2.5-7B & misinformation   & 961 & 0.39 & 0.521 & 0.800 & 0.859 \\
Mistral-7B & missing-evidence & 453 & 0.86 & 0.701 & 0.504 & 0.744 \\
Mistral-7B & misinformation   & 692 & 0.30 & 0.587 & 0.667 & 0.755 \\
\bottomrule
\end{tabular}
\end{table}

The pattern is clean and complementary. Semantic entropy detects
missing-evidence uncertainty (AUROC 0.67 to 0.70) but is near chance on
confident misinformation absorption (0.52 to 0.59): when the model
absorbs a false but coherent passage, its sampled answers agree on the
absorbed falsehood, so entropy is low even though the answer is wrong; this
is precisely the regime in which sampling-based entropy is known to be
ineffective \citep{ma2025semanticenergy}, and it aligns with the
epistemic-versus-aleatoric view of language-model uncertainty
\citep{ahdritz2024knowable}.
The retrieval descriptors are the mirror image, near chance on
missing-evidence (0.50 to 0.51, because retrieval scores do not reveal
the model's internal state) but strong on misinformation (0.67 to 0.80,
because contradictory and irrelevant passages have observable
signatures). The combined estimator captures both modes
(Table~\ref{tab:decomp}). This reconciles the whole paper. It explains
why the semantic entropy baseline collapses on injected conflict noise,
the headline of \S\ref{sec:baselines} and \S\ref{sec:generality}; why
semantic entropy looked competitive on purely natural noise
(\S\ref{sec:natnoise}), which is dominated by missing-evidence failures;
and why the combined estimator is the right design, since its two
signal families are complementary rather than redundant. The Qwen
missing-evidence cell has small $n = 51$ because that model abstains
heavily under irrelevance, so Mistral ($n = 453$) is the reliable
demonstration; the pattern is consistent across both generators.



\subsection{Purely natural retrieval noise}\label{sec:natnoise}
To rule out the objection that injected noise is synthetic, we ran the
pipeline with \emph{no} injection at all: the full SQuAD pool of about
2{,}000 paragraphs, Qwen2.5-7B, BM25 and dense retrieval pooled, 553
answered queries at a base hallucination rate of 0.212. The retriever's
own failures suffice to induce hallucination. When the gold paragraph is
retrieved, hallucination is 0.179 ($n = 507$); when the retriever misses
the gold paragraph it rises to 0.625 ($n = 16$); and closed-book sits at
0.533 ($n = 30$). A natural retrieval miss therefore raises hallucination
about 3.5 times with zero injected noise. On
this subtle natural noise the retrieval descriptors are weak and semantic
entropy is competitive (combined-minus-semantic-entropy $+0.017$, 95\% CI
$[-0.024, +0.059]$, not significant), the opposite of the injected
finding. The decomposition of \S\ref{sec:decomp} explains why: natural
retrieval noise is dominated by missing-evidence failures, the regime in
which semantic entropy works and the retrieval descriptors do not.



\section{Discussion}\label{sec:discussion}

\paragraph{Reconciling with the power of noise.}
\citet{cuconasu2024power} found that random unrelated documents can improve RAG
accuracy whereas high-ranked distractors hurt. Our irrelevance fillers are
topically plausible hard negatives, and their effect is monotonically harmful
(0.00 to 0.62 hallucination as severity rises): harm scales with the
plausibility of the noise, not its volume, and becomes silent exactly when
plausibility is highest.

\paragraph{Per-decision versus population certificates.}
C-RAG \citep{kang2024crag} certifies whether a RAG system is safe on average
under bounded drift; we show a system can satisfy such a bound while
per-decision calibration fails on every resplit under noise-regime selection.
The two are complementary: population certificates for system acceptance,
per-decision certification for runtime gating.

\paragraph{What the testbed buys and what it costs.}
The controlled world provides exact labels, causal attribution of
retrieval-induced error, and full coverage of a typed noise grid, none of which
naturalistic benchmarks offer; the cost is ecological validity (template text,
lexically detectable noise, BM25 retrieval). We regard it as a wind tunnel for
RAG, and the phenomena established here are mechanism-level claims, replicated
across two architectures, two content domains, six model families up to 9B, and
naturalistic SQuAD retrieval with two retrievers.

\paragraph{Limitations.}
Exact-match labels undercount paraphrase hallucinations, mitigated here by the
closed answer space; on the naturalistic data we additionally validate labels
against an independent large-model judge. Assumption~\ref{ass:shift}(ii) is an
idealisation; empirically the guarantee held across 200 adversarial resplits and
all sensitivity sweeps, the within-bin stress test of \S\ref{sec:selfdiag}
degraded it only to a $4\%$ violation rate, and the deployment monitor detects
the residual failure mode, but a shift that splits bins along an unmonitored
covariate could still degrade validity. The certification answer rate ($29\%$ at
$\alpha = 0.10$ under severe shift) is a substantial price; richer scores,
adaptive binning, and partial within-bin ordering are natural directions. Our
largest generator is 9B; the cross-architecture, cross-domain, and naturalistic
agreement address, but do not close, the frontier-scale question.

\appendix

\section{Reproducibility}\label{app:repro}
All corpora are generated by released scripts from fixed seeds; the released
dataset comprises 6{,}230 main-domain, 2{,}970 second-domain, and the
naturalistic generations, all with deterministic exact-match labels and a
pattern-based abstention class. Generators are public (flan-t5-base/small,
Qwen2.5-1.5B/3B/7B, Mistral-7B-v0.3, Gemma-2-9B, Phi-3.5-mini). The risk
estimator is gradient-boosted trees (150 trees, depth 3) under
question-grouped five-fold cross-validation, robust across logistic
regression, random forest, and a multilayer perceptron. The conformal
protocol uses 200 question-level resplits with clean-heavy calibration
($30\%$ noisy retention) and noise-heavy deployment ($50\%$ clean
retention), with $\alpha = 0.10$, $B = 10$, and $\delta = 0.10$ unless
swept; AUROC differences use a $2{,}000$-resample paired bootstrap over
questions. Complete code, data, cached out-of-fold scores, and the analysis
scripts that regenerate every table and figure are released at the
repository linked in the Declarations.

\backmatter

\bmhead{Declarations}
\paragraph{Funding.} The authors received no specific funding for this work.

\paragraph{Competing interests.} The authors declare no competing interests.

\paragraph{Author contributions.} Anmita Das and Shoumya Chowdhury
contributed equally to this work. A.D. and S.C. led the conceptualisation
and methodology, designed and implemented the experimental framework and
analysis pipeline, ran the experiments, prepared the figures, and wrote the
original draft. S.P. contributed to the conceptualisation and the framing of
the selective-generation problem, advised on the experimental design and the
statistical methodology, validated and interpreted the results, secured the
computational resources, administered the project, and revised the manuscript
through critical review and editing. All authors read and approved the final
manuscript.

\paragraph{Use of AI tools.} The authors used AI tools only for code
debugging and spell checking. All experimental designs, analyses, results,
and claims are the authors' own and were verified by the authors, who take
full accountability for the content of this article.

\paragraph{Data availability.} All datasets generated and analysed in this
study, together with the complete code for corpus construction, noise
injection, generation, and analysis, are publicly available at
\url{https://github.com/shoumya-chy/ACML-RAG}. The
study uses only publicly available, openly licensed models and synthetically
generated corpora; no human-subjects or proprietary data are used. The
naturalistic validation uses the publicly available SQuAD dataset together
with the openly licensed Qwen2.5-7B-Instruct and Mistral-7B-v0.3 generators.

\bibliography{refs}

\end{document}
